\section{Worked problems: reading a roofline}

\subsection{A 70-billion-parameter model on one H200}
\textbf{Problem.} A 70-billion-parameter model with Llama~3's shape (width
8,192, vocabulary 128,256, untied tables) is stored in FP8 on one H200. Predict
its batch-one decode rate from equation~\eqref{eq:decode-ceiling}, and list the
assumptions the prediction rests on.

\textbf{Solution.} At one byte per weight the file is about 70.5~GB. A decode
step reads every weight except the input table, which it gathers one row at a
time: $128{,}256 \times 8{,}192 = 1.05$~billion weights, 1.05~GB, so a step
reads about 69.5~GB. At the H200's 4.8~TB/s, $r_{\max} = 4.8/0.0695 = 69$
tokens per second. The prediction assumes (1) that the kernels reach the
datasheet bandwidth --- the measurements in this chapter suggest 70--90\% of it
is realistic, which gives 48--62 tokens per second; (2) that the KV cache is
small beside the weights, true for short conversations but not long ones (the
third problem); (3) that the FP8 scales add negligible bytes; and (4) that
nothing else --- kernel launches, sampling, synchronization --- adds time
between steps. The compute side needs no assumption at all: 139~GFLOP per token
at 69 tokens per second is 9.6~TFLOP/s, 0.5\% of the H200's dense FP8 peak.

\subsection{Where dequantization happens}
\textbf{Problem.} Weights are stored in four bits --- 4.5 bits with their
block scales, as in \texttt{Q4\_K} --- and multiplied in BF16 on an H100
(989~TFLOP/s, 3.35~TB/s). Derive $B^{*}$ from equation~\eqref{eq:batch-ridge}.
Then derive it again for an implementation that first expands the weights to
BF16 in a separate kernel, writes them to memory, and runs an ordinary BF16
matrix product on the result.

\textbf{Solution.} The ridge is $I^{*} = 989/3.35 = 295$~FLOP/byte. If the
matrix kernel reads the 4.5-bit weights and expands them in registers, a step
moves $b_w = 0.5625$ bytes per weight and does $2B$ FLOPs with it, so
$B^{*} = 0.5625 \times 295 / 2 = 83$. If dequantization is a separate kernel,
each weight costs 0.5625 bytes read, 2 bytes written in BF16 and 2 bytes read
back by the product: $b_w = 4.5625$, and $B^{*} = 4.5625 \times 295/2 = 673$.
Between 83 and 673 sequences the fused implementation is compute-bound and the
separate one is still waiting for memory; at batch one the separate one moves
eight times the bytes of the fused one and is eight times slower. The number of
bits in the file is not the number of bytes a step moves: that depends on where
the expansion happens, which is why every serious quantized kernel dequantizes
inside the matrix product (\Chref{weight-quantization}).

\subsection{Bytes that do not batch}
\textbf{Problem.} Qwen3-8B in BF16 decodes on an H100 at batch 1, 25 and 64,
with every sequence holding 512 or 4,096 tokens of 16-bit KV cache (144~KiB per
token). For each case compute the bytes and FLOPs of one step, the time
equation~\eqref{eq:step-bound} allows, the aggregate tokens per second, and which
roof binds. At what batch do the cache reads equal the weight reads?

\textbf{Solution.} A step reads 15.14~GB of weights once, whatever the batch,
plus $B \times S \times 147{,}456$ bytes of cache; it does $B \times
(15.14~\text{GFLOP} + 4 \times 36 \times 32 \times 128 \times S)$. At 512 tokens:
$B = 1$ moves 15.2~GB in at least 4.5~ms (220 tokens/s); $B = 64$ moves 20.0~GB
and does 0.99~TFLOP, 6.0~ms for memory against 1.0~ms for arithmetic, 10,700
tokens/s at an intensity of 49~FLOP/byte. At 4,096 tokens: $B = 1$ takes
4.7~ms (213 tokens/s); $B = 25$ moves 30.2~GB, 9.0~ms (2,770 tokens/s); $B = 64$
moves 53.8~GB --- more than three times the weights --- and needs 16.1~ms for
memory against 1.1~ms for arithmetic: 3,990 tokens/s at 21~FLOP/byte (all
derived). Memory binds in every case. The cache equals the weights at
$B = 15.14~\text{GB} / (S \times 147{,}456~\text{B})$: 25 sequences at 4,096
tokens, 200 at 512. Batching shares the weights, not the cache, so at long
context a large batch stays far down the bandwidth slope, and its throughput
grows with the batch only until the cache dominates the step. That is the
constraint \Chref{kv-cache} measures and Part~\ref{part:one-request}'s
cache-shrinking techniques relax.
